\documentclass[12pt,a4paper]{article}
\usepackage[utf8]{inputenc}
\usepackage[T1]{fontenc}
\usepackage[french]{babel}
\usepackage{amsmath, amssymb}
\usepackage{enumitem}
\usepackage[margin=2.5cm]{geometry}
\usepackage{hyperref}

\title{Résolution d’un système de congruences par le théorème des restes chinois}
\author{Nom: RANDRIAMANANTENA \\
Prénoms: Tsiky Ny Antsa \\
Matricule: 00550-80-23 \\
Parcours: Génie Logiciel}
\date{}
\begin{document}

\maketitle

\section*{Système à résoudre}
On cherche les entiers $x$ vérifiant :
\[
\begin{cases}
x \equiv 3 \pmod{7} \\
x \equiv 1 \pmod{11} \\
x \equiv 5 \pmod{13} \\
x \equiv 15 \pmod{17} \\
x \equiv 12 \pmod{23}
\end{cases}
\]
Les modules $7, 11, 13, 17, 23$ sont premiers entre eux deux à deux. Le théorème des restes chinois garantit alors une unique solution modulo $M = 7 \times 11 \times 13 \times 17 \times 23$.

\section*{1. Calcul du module global $M$}
\[
M = 7 \times 11 \times 13 \times 17 \times 23.
\]
Effectuons le produit pas à pas :
\begin{align*}
7 \times 11 &= 77,\\
77 \times 13 &= 1001,\\
1001 \times 17 &= 17017,\\
17017 \times 23 &= 391\,391.
\end{align*}
Donc
\[
\boxed{M = 391\,391}.
\]

\section*{2. Détermination des $M_i = M / m_i$}
Pour chaque congruence $x \equiv a_i \pmod{m_i}$, on pose $M_i = \dfrac{M}{m_i}$.

\begin{align*}
M_1 &= \frac{391\,391}{7} = 55\,913, \\
M_2 &= \frac{391\,391}{11} = 35\,581, \\
M_3 &= \frac{391\,391}{13} = 30\,107, \\
M_4 &= \frac{391\,391}{17} = 23\,023, \\
M_5 &= \frac{391\,391}{23} = 17\,017.
\end{align*}

\section*{3. Calcul des inverses modulaires $y_i$}
Pour chaque $i$, on cherche $y_i$ tel que $M_i \cdot y_i \equiv 1 \pmod{m_i}$.

\subsection*{Pour $i=1$ : $M_1 \equiv 55\,913 \pmod{7}$}
\[
55\,913 \div 7 = 7\,987 \times 7 = 55\,909,\quad \text{reste } 4.
\]
Donc $M_1 \equiv 4 \pmod{7}$.
On résout $4y \equiv 1 \pmod{7}$.
\[
4 \times 2 = 8 \equiv 1 \pmod{7} \quad\Longrightarrow\quad y_1 = 2.
\]

\subsection*{Pour $i=2$ : $M_2 \equiv 35\,581 \pmod{11}$}
\[
35\,581 \div 11 = 3\,234 \times 11 = 35\,574,\quad \text{reste } 7.
\]
Donc $M_2 \equiv 7 \pmod{11}$.
Résolvons $7y \equiv 1 \pmod{11}$.
\[
7 \times 8 = 56 = 5\times 11 + 1 \equiv 1 \pmod{11} \quad\Longrightarrow\quad y_2 = 8.
\]

\subsection*{Pour $i=3$ : $M_3 \equiv 30\,107 \pmod{13}$}
\[
30\,107 \div 13 = 2\,315 \times 13 = 30\,095,\quad \text{reste } 12 \equiv -1 \pmod{13}.
\]
Donc $M_3 \equiv -1 \pmod{13}$.
L’inverse de $-1$ modulo $13$ est $-1$ lui-même car $(-1)\times(-1)=1$.
Ainsi $y_3 = -1 \equiv 12 \pmod{13}$ ; on prendra $y_3 = 12$.

\subsection*{Pour $i=4$ : $M_4 \equiv 23\,023 \pmod{17}$}
\[
23\,023 \div 17 = 1\,354 \times 17 = 23\,018,\quad \text{reste } 5.
\]
Donc $M_4 \equiv 5 \pmod{17}$.
On résout $5y \equiv 1 \pmod{17}$.
\[
5 \times 7 = 35 = 2\times 17 + 1 \equiv 1 \pmod{17} \quad\Longrightarrow\quad y_4 = 7.
\]

\subsection*{Pour $i=5$ : $M_5 \equiv 17\,017 \pmod{23}$}
\[
17\,017 \div 23 = 739 \times 23 = 16\,997,\quad \text{reste } 20 \equiv -3 \pmod{23}.
\]
Donc $M_5 \equiv -3 \pmod{23}$.
On cherche $y$ tel que $(-3)y \equiv 1 \pmod{23}$, soit $3y \equiv -1 \equiv 22 \pmod{23}$.
\[
3 \times 15 = 45 = 1\times 23 + 22 \equiv 22 \pmod{23}.
\]
Donc $y_5 = 15$ (car $3\times15=45 \equiv 22 \pmod{23}$).
(Vérification : $(-3)\times15 = -45 = -2\times23 + 1 \equiv 1 \pmod{23}$.)

\section*{4. Construction de la solution}
La solution est donnée par
\[
x = \sum_{i=1}^{5} a_i \, M_i \, y_i \pmod{M}.
\]

Calculons chaque terme :
\begin{align*}
a_1 M_1 y_1 &= 3 \times 55\,913 \times 2 = 335\,478, \\
a_2 M_2 y_2 &= 1 \times 35\,581 \times 8 = 284\,648, \\
a_3 M_3 y_3 &= 5 \times 30\,107 \times 12 = 1\,806\,420, \\
a_4 M_4 y_4 &= 15 \times 23\,023 \times 7 = 2\,417\,415, \\
a_5 M_5 y_5 &= 12 \times 17\,017 \times 15 = 3\,063\,060.
\end{align*}

Somme :
\[
S = 335\,478 + 284\,648 + 1\,806\,420 + 2\,417\,415 + 3\,063\,060 = 7\,907\,021.
\]

\section*{5. Réduction modulo $M$}
On calcule le reste de $S$ modulo $391\,391$ :
\[
7\,907\,021 \div 391\,391 \approx 20.2,\quad
391\,391 \times 20 = 7\,827\,820.
\]
\[
\text{Reste} = 7\,907\,021 - 7\,827\,820 = 79\,201.
\]
Ainsi
\[
\boxed{x \equiv 79\,201 \pmod{391\,391}}.
\]

La plus petite solution positive est $79\,201$.

\section*{6. Vérification}
On vérifie chaque congruence :
\begin{align*}
79\,201 \bmod 7   &= 3, \quad (\text{car } 7 \times 11\,314 = 79\,198,\;\; 79\,201-79\,198 = 3),\\
79\,201 \bmod 11  &= 1, \quad (11 \times 7\,200 = 79\,200,\;\; \text{reste }1),\\
79\,201 \bmod 13  &= 5, \quad (13 \times 6\,092 = 79\,196,\;\; \text{reste }5),\\
79\,201 \bmod 17  &= 15, \quad (17 \times 4\,658 = 79\,186,\;\; \text{reste }15),\\
79\,201 \bmod 23  &= 12, \quad (23 \times 3\,443 = 79\,189,\;\; 79\,201-79\,189 = 12).
\end{align*}
Toutes les conditions sont satisfaites.

\section*{7. Ensemble des solutions}
L’ensemble complet des entiers $x$ solutions est
\[
\{ \, 79\,201 + 391\,391\,k \mid k \in \mathbb{Z} \, \}.
\]

\end{document} 
